\chapter{\label{ch:3-methods}Search}

\minitoc

\section{Introduction}
The task of obtaining the stabiliser decomposition for classical simulation purposes
has only been discussed at a high level so far.
While we have the various methods of magic state, cat state decompositions, 
along with grap cut decompositions,
we have not yet detailed the specific steps taken to obtain the stabiliser decompositions
in the most efficient way so far.

In modeling our task, we can think of the stabiliser decomposition as a search problem, 
where we are searching for the most efficient stabiliser decomposition.
We can visualize an example as in \autoref{fig:amplitude}, where $V$ is a Clifford+T diagram at the root.
The directed edges represent the decomposition used from the parent to the child node, 
and the leaves $d_i$ are the possible stabiliser decompositions.

\clearpage

\begin{figure}[h!]
\tikzfig{search-tree}
\caption{Example search tree for a stabiliser decomposition}
\end{figure}

To traverse the search tree, we might take the following steps \cite{Kissinger2022cat,Kissinger_2022}.

\begin{algorithm}
\label{alg:stabiliser-decomposition}
\begin{enumerate}    
    \item Perform necessary spider $\SpiderRule$usions to 
    \hyperref[prop:zx-to-graph]{reduce to graph-like form}.
    \item Apply a chosen decomposition ("choose an edge").
    \item Simplify the resulting Clifford+T diagram using \hyperref[alg:zx-simplify]{ZX-simplify}.
    Some Clifford spiders may be removed, and some $T$-like spiders may be combined.
    \item Repeat steps 1-3 until the $T$-count is 0.
\end{enumerate}
\end{algorithm}

The complexity from using this algorithm was found to be $O(N^3 + 2^{\alpha t})$, 
where $N$ is the number of gates in the original circuit.
Clearly, any circuit with a $T$-count that is not insignificant will mean that 
the complexity is dominated by $O(2^{\alpha t})$.

In this thesis, we will focus on the second step of the algorithm, 
where we choose the decomposition to apply at each node.
There are a variety of ways to choose a decomposition, detailed in the next few sections.

\section{Greedy Algorithm}

The most obvious way to obtain an efficient decomposition is to, at each step, use the decomposition 
that has the smallest associated value of $\alpha$ as possible.
This is essentially a greedy algorithm, and is the method used in \cite{Kissinger2022cat,Kissinger_2022}.
Intuitively, this should give us a decomposition with an $\alpha$ upper bounded by 
the worst decomposition available.
More formally, we have the following lemma.

\begin{lemma}
    Let $d$ be a stabiliser decomposition, 
    and let $\alpha_1,...,\alpha_n$ and $t_1,...,t_n$ be the $\alpha$ values and $T$-count reductions
    of each of the $n$ decompositions to obtain $d$ from the initial Clifford+T diagram.
    Then the efficiency $\alpha$ value of the decomposition $d$ is given by
\begin{equation}
    \alpha \leq \frac{\sum_{i=1}^n \alpha_i t_i}{\sum_{i=1}^n t_i}
\end{equation}
\end{lemma}

\begin{proof}
    Each decomposition gives $\leq 2^{\alpha_i t_i}$ terms, for a total of
    \[\leq \prod_{i=1}^{n} 2^{\alpha_i t_i} = 2^{\sum_{i=1}^n \alpha_i t_i}\]
    terms. Since the final $T$-count is 0, the initial $T$-count is $\sum_{i=1}^n t_i$.
    By definition, we have 
    \[\alpha \leq \frac{\sum_{i=1}^n \alpha_i t_i}{\sum_{i=1}^n t_i}\]
\end{proof}

This tells us that the final $\alpha$ value is essentially a weighted average of the $\alpha_i$ values,
weighted by their $T$-count reductions.
The actual $\alpha$ value in practice is likely to be much lower than this upper bound,
because of step 3 in the \autoref{alg:stabiliser-decomposition}, 
where we simplify the diagram using \nameref{alg:zx-simplify}.
In \cite{Kissinger_2022}, although the BSS decomposition has the $\alpha \approx 0.468$,
the actual $\alpha$ value obtained in experiments was lower, at $\approx 0.4$ \cite{sutcliffe2024proccut}.

\section{Graph Structure Heuristic}
\label{sec:graph-cut}

Another way to choose a decomposition is to use the graph structure of the Clifford+T diagram.
A disadvantage of the greedy algorithm is that it does not take into account the structure of the diagram.
When we only consider the $\alpha$ value of the decomposition, 
we fail to take into account the \nameref{alg:zx-simplify} step of Algorithm \ref{alg:stabiliser-decomposition}.
Since this step can combine $T$-like spiders,
it means that decompositions with seemingly higher $\alpha$ values may actually be more efficient.

In \cite{sutcliffe2024proccut}, the heuristic used considers the structure of CNOT gates
placed in between $T$ gates.
The following observation was made.

\begin{equation}
    \tikzfig{picom_01}
\end{equation}

Using the graph cut (\ref{eq:spider_cut}), 
which without simplification has $\alpha = 1$ if the spider being cut is a $T$-like spider,
we can instead obtain a decomposition with $\alpha \approx 0.333$, with simplication,
for the reduction of 3 $T$-like spiders at the cost of 2 terms.
In fact, this structure can be stacked in the following way.

\begin{equation}
    \tikzfig{picom_grouped}
\end{equation}